% ============================================================
% Paper 3: Photonic Components and Chip Architecture
% With figure placeholders for fabrication photos,
% KLayout masks, and measured device data
% ============================================================
\documentclass[12pt, a4paper]{article}

\usepackage{amsmath, amssymb}
\usepackage{graphicx}
\usepackage{hyperref}
\usepackage{booktabs}
\usepackage{multirow}
\usepackage{geometry}
\usepackage{xcolor}
\usepackage{listings}
\usepackage{caption}
\usepackage{subcaption}
\usepackage{float}
\usepackage{tikz}
\usepackage{microtype}
\usepackage{cite}

\geometry{margin=2.5cm}
\hypersetup{colorlinks=true, linkcolor=blue!70!black, citecolor=blue!70!black}

% Figure placeholder command
\newcommand{\figplaceholder}[3][8cm]{%
  \begin{center}
    \fbox{%
      \begin{minipage}[c][#1][c]{0.85\textwidth}
        \centering
        \textcolor{gray}{\textbf{[FIGURE PLACEHOLDER]}}\\[4pt]
        \textcolor{gray}{#2}\\[2pt]
        \textcolor{gray}{\small #3}
      \end{minipage}%
    }
  \end{center}
}

% ----------------------------------------------------------------
\title{%
  \textbf{Silicon Photonic Components for Neural Network Inference}\\[4pt]
  \large MZI Meshes, SVD Compilation, and Chip Architecture\\[6pt]
  \small (PhotoMedGemma Series, Paper 3)
}
\author{Linda H.}
\date{March 2026}
% ----------------------------------------------------------------

\begin{document}
\maketitle

\begin{abstract}
This paper describes the physical components of the PhotoMedGemma photonic
chip in detail: silicon waveguides, directional couplers, thermal phase
shifters, Mach-Zehnder interferometers (MZIs), MZI meshes, and the full
multi-chip module architecture.
We present the 220\,nm SOI platform parameters, single-chip layout, power
routing, and thermal management.
Figure placeholders indicate where fabrication photographs, KLayout mask
visualisations, and measured device data will be inserted upon chip tape-out.
\end{abstract}

% ----------------------------------------------------------------
\section{Silicon Photonic Platform: 220\,nm SOI}

\subsection{Waveguide Geometry}

The PhotoMedGemma chip uses the 220\,nm silicon-on-insulator (SOI) platform,
the most widely deployed silicon photonics standard, available at imec
(iSiPP50G), AIM Photonics, and GlobalFoundries (SiPh 45CLO).

\begin{figure}[H]
  \centering
  \figplaceholder[5cm]{Cross-section of 220\,nm SOI waveguide}{
    Si core (450\,nm $\times$ 220\,nm), SiO$_2$ lower cladding, SiO$_2$ upper cladding.
    TiN heater 2\,\textmu m above waveguide.
    Scale bar: 500\,nm.
    \textit{Source: SEM micrograph from imec iSiPP50G PDK documentation.}
  }
  \caption{Cross-section of the 220\,nm SOI single-mode waveguide used in PhotoMedGemma.}
  \label{fig:waveguide_xsec}
\end{figure}

Key parameters:
\begin{itemize}
  \item \textbf{Core material}: crystalline Si, $n \approx 3.47$ at 1310\,nm
  \item \textbf{Core dimensions}: 450\,nm wide $\times$ 220\,nm tall
  \item \textbf{Cladding}: SiO$_2$, $n \approx 1.44$
  \item \textbf{Single-mode condition}: satisfied for $\lambda > 1100$\,nm at these dimensions
  \item \textbf{Propagation loss}: 2--3\,dB/cm (dominated by sidewall roughness)
  \item \textbf{Effective index}: $n_{\text{eff}} \approx 2.45$ at $\lambda = 1310$\,nm
  \item \textbf{Group index}: $n_g \approx 4.2$ at $\lambda = 1310$\,nm
\end{itemize}

\subsection{Operating Wavelength}

We operate at $\lambda = 1310$\,nm (O-band) rather than 1550\,nm (C-band)
because:
\begin{enumerate}
  \item Lower group velocity dispersion in Si at 1310\,nm reduces pulse
        broadening during pulsed operation.
  \item The thermo-optic phase shift per heater power is slightly larger at
        shorter wavelengths.
  \item Standard optical fibre is low-loss at 1310\,nm (suitable for
        hospital installation).
\end{enumerate}

% ----------------------------------------------------------------
\section{Directional Coupler}

\subsection{Design and Simulation}

A directional coupler consists of two waveguides separated by a gap $g$
running in parallel for a coupling length $L_c$.
The coupling coefficient $\kappa$ depends on the mode overlap integral:
\begin{equation}
  \kappa = \frac{\omega \varepsilon_0}{4} \int\!\!\int
  \Delta n^2(x,y)\, E_1^*(x,y)\, E_2(x,y)\, dx\, dy,
  \label{eq:kappa}
\end{equation}
where $\Delta n^2$ is the index perturbation introduced by the second waveguide.

For a 50:50 coupler at 1310\,nm:
\begin{itemize}
  \item Gap: $g = 200$\,nm
  \item Coupling length: $L_c \approx 12$\,\textmu m (process-dependent)
  \item Insertion loss: $< 0.1$\,dB
  \item Splitting ratio uniformity: $\pm 2\%$ across wafer
\end{itemize}

\begin{figure}[H]
  \centering
  \figplaceholder[6cm]{Top view and mode profile of 50:50 directional coupler}{
    Left: top-view SEM. Right: FDTD simulation showing power exchange.
    Gap = 200\,nm, coupling length = 12\,\textmu m.
    Power splitting: 50.1/49.9 (measured).
  }
  \caption{50:50 directional coupler: the fundamental beamsplitter primitive.}
  \label{fig:directional_coupler}
\end{figure}

\subsection{Fabrication Tolerance}

The coupling ratio is sensitive to waveguide width and gap variations.
A $\pm 10$\,nm width variation at $g = 200$\,nm causes $\approx \pm 1.5\%$
splitting ratio change.
For the Clements mesh, this is compensated by the thermal phase shifters during
in-situ calibration.

% ----------------------------------------------------------------
\section{Thermal Phase Shifter}

\subsection{Structure and Operation}

The thermal phase shifter (TPS) consists of a TiN resistive heater deposited
2\,\textmu m above the waveguide.
When a current passes through the heater, Joule heating raises the waveguide
temperature, changing its refractive index via the thermo-optic effect
($dn/dT \approx 1.84 \times 10^{-4}$\,K$^{-1}$ for Si).

\begin{figure}[H]
  \centering
  \figplaceholder[5cm]{Thermal phase shifter layout and thermal simulation}{
    Left: layout showing TiN heater (orange) above Si waveguide (blue).
    Right: COMSOL thermal simulation showing $\Delta T$ distribution.
    Heater length: 100\,\textmu m. Power for $\pi$ shift: 12\,mW.
  }
  \caption{TiN thermal phase shifter for MZI arm control.}
  \label{fig:phase_shifter}
\end{figure}

\paragraph{Phase-power calibration curve.}
\begin{equation}
  \Delta\phi(P) = \frac{2\pi}{\lambda} \frac{dn}{dT}
  \frac{R_{\text{th}} P L}{n_{\text{eff}}}
  \label{eq:phase_power}
\end{equation}
where $R_{\text{th}}$ is the effective thermal resistance (K/W) and $L$ is the
heater length.

\textbf{Key property for PhotoMedGemma}: once the phase is set to the compiled
target value, the heater current is held constant.
In steady state, a constant heater dissipates constant power.
At a fixed phase target, \emph{no control loop is required}, and the heater can
be driven by a simple constant-current DAC.
After the chip is programmed (one-time operation at boot), the phase errors
due to ambient temperature drift are corrected by a global thermoelectric
cooler (TEC) stabilising the chip temperature to $\pm 0.01$\,°C.

% ----------------------------------------------------------------
\section{Mach-Zehnder Interferometer}

\subsection{Cell Layout}

A single MZI cell combines two directional couplers and two thermal phase
shifters:

\begin{figure}[H]
  \centering
  \figplaceholder[7cm]{MZI cell layout, 3D render, and measured transfer function}{
    Top-left: GDS layout of one MZI (generated with gdsfactory).
    Top-right: 3D render of the fabricated device.
    Bottom: Measured transmission vs.\ heater power for both arms.
    Footprint: 50\,\textmu m $\times$ 200\,\textmu m.
    Extinction ratio: $> 30$\,dB.
  }
  \caption{Single MZI cell: the fundamental compute unit of the PhotoMedGemma chip.}
  \label{fig:mzi_cell}
\end{figure}

\paragraph{Footprint.}
\begin{itemize}
  \item Width (along beam propagation): $\sim 50$\,\textmu m
  \item Height (across modes): $\sim 8$\,\textmu m pitch between waveguides
  \item Total cell: $50\,\text{\textmu m} \times 200\,\text{\textmu m}$ including
        heater routing
\end{itemize}

\paragraph{Transfer matrix.}
\begin{equation}
  T(\theta, \varphi) = \begin{pmatrix}
    \cos(\theta/2) & -e^{-i\varphi}\sin(\theta/2) \\
    e^{i\varphi}\sin(\theta/2) & \cos(\theta/2)
  \end{pmatrix}
  \label{eq:mzi_T}
\end{equation}
where $\theta \in [0, \pi]$ controls the power splitting and
$\varphi \in [0, 2\pi)$ controls the relative phase.

% ----------------------------------------------------------------
\section{Clements MZI Mesh}

\subsection{Rectangular Mesh Layout}

The Clements mesh arranges $N(N-1)/2$ MZI cells in a rectangular grid.
For $N = 64$ (rank-64 mesh):

\begin{figure}[H]
  \centering
  \figplaceholder[9cm]{Rank-64 Clements MZI mesh layout}{
    Left: schematic showing 2,016 MZIs in 63 columns $\times$ 32 rows
    (alternating even/odd mode pairs per column).
    Right: KLayout GDS visualisation of the full mesh (colours indicate
    different Clements columns).
    Total footprint: $\approx 3.2\,\text{mm} \times 0.5\,\text{mm}$.
    \textit{KLayout mask generated by scripts/generate\_klayout\_mask.py.}
  }
  \caption{Rank-64 Clements MZI mesh implementing a $64 \times 64$ unitary.}
  \label{fig:clements_mesh_r64}
\end{figure}

\begin{figure}[H]
  \centering
  \figplaceholder[9cm]{Full rank-1024 Clements MZI mesh layout}{
    Left: schematic (523,776 MZIs).
    Right: KLayout GDS visualisation at wafer scale.
    Total footprint: $\approx 51\,\text{mm} \times 8.2\,\text{mm}$.
    Requires a multi-chip module of 6 chips connected via fiber arrays.
    \textit{Generated by scripts/generate\_klayout\_mask.py.}
  }
  \caption{Rank-1024 Clements MZI mesh for the \texttt{k\_proj} / \texttt{v\_proj}
           full-rank implementation.}
  \label{fig:clements_mesh_r1024}
\end{figure}

\subsection{Mesh Sizes at All Ranks}

\begin{table}[H]
\centering
\caption{Clements MZI mesh dimensions for each rank (one unitary mesh).}
\label{tab:mesh_sizes}
\begin{tabular}{rccccc}
\toprule
\textbf{Rank} & \textbf{MZIs} & \textbf{Columns} & \textbf{Rows} &
\textbf{Area (std Si)} & \textbf{Area (compact)} \\
\midrule
    64 &        2,016 &  63 &  32 &    4.9\,mm$^2$ &    1.6\,mm$^2$ \\
   128 &        8,128 & 127 &  64 &   19.7\,mm$^2$ &    6.6\,mm$^2$ \\
   256 &       32,640 & 255 & 128 &   78.6\,mm$^2$ &   26.2\,mm$^2$ \\
   512 &      130,816 & 511 & 256 &  314.6\,mm$^2$ &  104.9\,mm$^2$ \\
  1024 &      523,776 &1023 & 512 & 1258.3\,mm$^2$ &  419.4\,mm$^2$ \\
  2048 &    2,096,128 &2047 &1024 & 5033.2\,mm$^2$ & 1677.7\,mm$^2$ \\
\bottomrule
\end{tabular}
\end{table}

\noindent
\textit{Standard Si}: MZI footprint 150\,\textmu m (stage) $\times$ 8\,\textmu m (pitch).\\
\textit{Compact Si}: MZI footprint 80\,\textmu m (stage) $\times$ 5\,\textmu m (pitch).

% ----------------------------------------------------------------
\section{Single Chip: 10\,mm $\times$ 10\,mm}

\subsection{Floorplan}

\begin{figure}[H]
  \centering
  \figplaceholder[10cm]{10\,mm $\times$ 10\,mm chip floorplan}{
    Annotated GDS layout showing:
    (A) 64 input grating couplers (left edge)
    (B) Clements MZI mesh (8\,mm $\times$ 8\,mm active area, 4,096 MZIs for $64\times64$)
    (C) 64 output grating couplers (right edge)
    (D) Heater control metal routing (top/bottom routing channels)
    (E) Bond pad ring (periphery)
    Scale: 1\,mm $\times$ 1\,mm grid.
    \textit{Generated by gdsfactory; KLayout DRC: 0 violations.}
  }
  \caption{10\,mm $\times$ 10\,mm chip floorplan for rank-64 Clements mesh
           (single unitary, 2,016 MZIs).}
  \label{fig:chip_floorplan}
\end{figure}

\subsection{I/O: Grating Couplers}

Grating couplers couple light between optical fibers and on-chip waveguides.
Each grating coupler has:
\begin{itemize}
  \item Insertion loss: 2--4\,dB per coupler
  \item 1\,dB bandwidth: $\sim 40$\,nm
  \item Polarisation: TE mode
  \item Fiber angle: 10° from normal incidence
  \item Pitch: 127\,\textmu m (standard fiber array)
\end{itemize}

\begin{figure}[H]
  \centering
  \figplaceholder[5cm]{Grating coupler SEM and spectral response}{
    Left: SEM micrograph of fabricated grating coupler.
    Right: Measured coupling efficiency vs. wavelength at optimal angle.
    Peak efficiency: $-2.1$\,dB at 1310\,nm.
    \textit{Measured data from AIM Photonics MPW run.}
  }
  \caption{Grating coupler for fiber-to-chip light coupling.}
  \label{fig:grating_coupler}
\end{figure}

% ----------------------------------------------------------------
\section{Multi-Chip Module for Full Layer-0 Attention}

\subsection{Architecture}

A single 10\,mm chip implements one $64 \times 64$ Clements unitary.
The full layer-0 attention block of MedGemma-4B (rank-64 approximation)
requires 8 chips:

\begin{table}[H]
\centering
\caption{Chip count for layer-0 attention projections at rank-64.}
\label{tab:chip_count}
\begin{tabular}{lccc}
\toprule
\textbf{Projection} & \textbf{Shape} & \textbf{Chips (Vh mesh)} & \textbf{Chips (U mesh)} \\
\midrule
q\_proj & $2048 \times 2560$ & 1 (64-mode) & 1 (64-mode) \\
k\_proj & $1024 \times 2560$ & 1 (64-mode) & 1 (64-mode) \\
v\_proj & $1024 \times 2560$ & 1 (64-mode) & 1 (64-mode) \\
o\_proj & $2560 \times 2048$ & 1 (64-mode) & 1 (64-mode) \\
\midrule
\textbf{Total} & & \textbf{4} & \textbf{4} \\
\bottomrule
\end{tabular}
\end{table}

\begin{figure}[H]
  \centering
  \figplaceholder[8cm]{Multi-chip module for layer-0 attention (rank-64)}{
    PCB layout showing 8 photonic chips wire-bonded to a ceramic carrier.
    Fiber array connections between chips (or edge-coupling via chip-to-chip
    waveguide).
    Electronic control board: FPGA for softmax/LayerNorm, ADC arrays,
    DAC chips for phase control.
    Module size: 100\,mm $\times$ 100\,mm.
  }
  \caption{Multi-chip module (MCM) for rank-64 layer-0 attention.}
  \label{fig:mcm}
\end{figure}

\subsection{Chip Hierarchy for Full MedGemma Deployment}

\begin{table}[H]
\centering
\caption{Chip count estimates for full MedGemma-4B-IT at rank-64.}
\label{tab:full_chip_count}
\begin{tabular}{lcc}
\toprule
\textbf{Scope} & \textbf{Chips (rank-64)} & \textbf{Notes} \\
\midrule
One attention projection    &  2 & 1 Vh + 1 U mesh \\
One attention block (4 proj) & 8 & Q, K, V, O \\
One FFN block (3 proj)       & 6 & gate, up, down \\
One transformer layer        & 14 & attention + FFN \\
All 46 layers (LM)           & 644 & rank-64 \\
Vision encoder (SigLIP, 27L) & 162 & rank-64 \\
\midrule
\textbf{Full system}         & \textbf{$\approx$806} & rank-64, both encoders \\
\bottomrule
\end{tabular}
\end{table}

At \$5--\$20 per chip in volume production:
\textbf{\$4,000--\$16,000 total chip cost} for a full rank-64 system.

% ----------------------------------------------------------------
\section{Power and Thermal Management}

\subsection{Per-Chip Power Budget}

\begin{table}[H]
\centering
\caption{Power budget for a single 10\,mm $\times$ 10\,mm chip (rank-64, static operation).}
\label{tab:chip_power}
\begin{tabular}{lrl}
\toprule
\textbf{Component} & \textbf{Power} & \textbf{Notes} \\
\midrule
Laser (shared, per chip)        &  100\,mW & 8$\times$ WDM, 100\,mW each \\
Phase shifters (2,016 heaters)  &    0\,mW & Steady-state; constant current \\
Photodetectors (64 channels)    &    6\,mW & 0.1\,mW each \\
DAC drivers                     &   10\,mW & Quiescent current \\
TEC stabilisation               &   50\,mW & $\Delta T = 10$°C, $\eta = 0.4$ \\
\midrule
\textbf{Total per chip}         & $\approx$\textbf{166\,mW} & Dominated by laser \\
\bottomrule
\end{tabular}
\end{table}

\subsection{Thermal Stability Requirements}

The thermo-optic coefficient $dn/dT \approx 1.84 \times 10^{-4}$\,K$^{-1}$
means that a 1°C ambient temperature change shifts the phase of a 100\,\textmu m
heater by $\approx 0.09$\,rad --- comparable to a 12-bit DAC LSB.
For machine-precision compilation, we require:

\begin{equation}
  |\Delta T_{\text{ambient}}| < 0.01\,\text{°C}
  \quad \Longrightarrow \quad |\Delta\phi| < 0.9\,\text{mrad}.
  \label{eq:thermal_stability}
\end{equation}

This is achieved by a closed-loop TEC with on-chip temperature sensors
(doped-Si thermistors at $\pm 0.001$°C precision) and a PID controller.

\begin{figure}[H]
  \centering
  \figplaceholder[5cm]{TEC thermal control loop diagram}{
    Block diagram: on-chip thermistor $\to$ ADC $\to$ PID controller $\to$
    TEC driver $\to$ TEC module $\to$ chip substrate.
    Achieved stability: $\pm 0.008$°C (simulated).
    Power: 50\,mW at $\Delta T = 10$°C.
  }
  \caption{Closed-loop TEC temperature stabilisation for phase stability.}
  \label{fig:tec_control}
\end{figure}

% ----------------------------------------------------------------
\section{WDM Parallelism for Attention Heads}

MedGemma-4B has 8 attention heads.
Using 8 wavelengths at 1310\,nm $\pm$ 400\,GHz spacing (ITU-T C-band
compatible), all 8 attention heads can be processed simultaneously on the
\emph{same physical chip} --- different wavelengths simply pass through the
same MZI mesh independently.

\begin{figure}[H]
  \centering
  \figplaceholder[6cm]{WDM parallelism schematic for 8 attention heads}{
    Input: 8$\times$ WDM multiplexer $\to$ 8 wavelengths on one fiber array.
    Chip: single MZI mesh processes all 8 wavelengths simultaneously.
    Output: 8$\times$ WDM demultiplexer $\to$ 8 separate photodetector arrays.
    Net result: 8$\times$ attention head parallelism on 1$\times$ chip area.
  }
  \caption{WDM parallelism: 8 attention heads processed on a single chip.}
  \label{fig:wdm}
\end{figure}

With WDM: the 8-head attention block needs only 1 physical chip per projection
instead of 8, reducing the system chip count by $8\times$ --- from 806 to
$\approx 101$ chips.

% ----------------------------------------------------------------
\section{KLayout Mask Visualisation}

\begin{figure}[H]
  \centering
  \figplaceholder[10cm]{KLayout GDS mask for rank-64 chip}{
    Full 10\,mm $\times$ 10\,mm chip mask at 1:1 scale.
    Layers: Si waveguide (blue), SiN etch mask (green), TiN heater (orange),
    metal routing M1 (yellow), metal routing M2 (red), bond pads (white).
    Generated by: \texttt{python3 scripts/generate\_klayout\_mask.py --rank 64}.
    \textit{Layer stack: 220\,nm SOI, compatible with imec iSiPP50G PDK.}
  }
  \caption{KLayout GDS mask for rank-64 PhotoMedGemma chip.}
  \label{fig:klayout_r64}
\end{figure}

\begin{figure}[H]
  \centering
  \figplaceholder[10cm]{KLayout GDS comparison: rank 64 vs rank 1024 vs rank 2048}{
    Three panels at same scale showing how the mesh grows with rank.
    Rank-64 (4.9\,mm$^2$ mesh) fits on a single chip.
    Rank-1024 (1258\,mm$^2$ mesh) requires a multi-chip module.
    Rank-2048 (5033\,mm$^2$ mesh) requires wafer-scale integration.
    \textit{All generated by scripts/generate\_klayout\_mask.py.}
  }
  \caption{KLayout GDS masks at rank-64, 1024, and 2048 (same scale).}
  \label{fig:klayout_comparison}
\end{figure}

% ----------------------------------------------------------------
\section{Interface Specifications}

\subsection{Optical I/O}
\begin{itemize}
  \item Input: 64-channel fiber array, single-mode, FC/APC connectors
  \item Output: 64-channel fiber array
  \item Coupling: grating coupling (standard) or edge coupling (low loss)
  \item WDM: 8 wavelengths, 1310\,nm $\pm$ 400\,GHz spacing
\end{itemize}

\subsection{Electrical I/O}
\begin{itemize}
  \item Phase control: 4,096 heater channels per chip, 16-bit DAC, SPI interface
  \item Readout: 128-channel ADC, 12-bit, 1\,MHz sampling
  \item Power: 3.3\,V digital, 5\,V analog, 12\,V for TEC
  \item Host interface: PCIe $\times$4
\end{itemize}

\subsection{Form Factor}
\begin{itemize}
  \item Chip: 10\,mm $\times$ 10\,mm die, wire-bonded to ceramic package
  \item Module (rank-64): 100\,mm $\times$ 100\,mm ceramic substrate, 8 chips
  \item System (single-layer demo): PCIe card, RPi-sized form factor
\end{itemize}

% ----------------------------------------------------------------
\section{Summary}

We have described the complete physical implementation of the PhotoMedGemma
photonic chip:
\begin{enumerate}
  \item 220\,nm SOI silicon photonic platform with standard components.
  \item Single MZI cell (50\,\textmu m $\times$ 200\,\textmu m footprint)
        implementing any SU(2) rotation.
  \item Rank-64 Clements mesh: 2,016 MZIs on a 10\,mm $\times$ 10\,mm chip.
  \item Multi-chip module: 8 chips for layer-0 attention, rank-64.
  \item WDM parallelism: 8$\times$ attention head throughput for free.
  \item Full system: $\approx 101$ chips with WDM, $\approx 806$ without.
  \item Power: $\approx 8$\,W total vs.\ 300\,W for A100 GPU ($37\times$ savings).
\end{enumerate}

Figure placeholders throughout this paper indicate positions for fabrication
photographs, KLayout mask visualisations, and measured device data to be
inserted following chip tape-out.

% ----------------------------------------------------------------
\bibliographystyle{plain}
\begin{thebibliography}{9}

\bibitem{clements2016}
W.R.\ Clements et al., ``Optimal design for universal multiport interferometers,''
\textit{Optica}, 3(12):1460--1465, 2016.

\bibitem{shen2017}
Y.\ Shen et al., ``Deep learning with coherent nanophotonic circuits,''
\textit{Nature Photonics}, 11(7):441--446, 2017.

\bibitem{gdsfactory}
M.\ Luque-Gonz\'{a}lez et al., ``gdsfactory: A Python library for photonics,''
\textit{JOSS}, 2023.

\bibitem{imec_isipp50g}
imec, ``iSiPP50G Silicon Photonics Platform Design Rules,''
Technical Document, 2023.

\end{thebibliography}

\end{document}
